\paragraph{检测\trQuasiRepresent{}元。}\label{sub:detecting-quasi-representatives}
给定一个 \trACategory{} $\scr{A}$ 的\trObject{} $Y_0$、$Y_1$、$Y_2$，若能有一个
显式的判据来检测 $Y_2$ 何时可以取作 $\scr{T}_{Y_0}\scr{Y}_1$ 的%
\trQuasiRepresent{}元，那会很有用。人们可以把这件事建立在%
\cref{thm:quasi-represents-characterization}之上，该引理给出了定义\trMorphism{}%
$\scr{Y}_2\to\scr{T}_{Y_0}\scr{Y}_1$ 所需的数据；然而我们不愿那样做，而是宁愿
着眼于一个方向相反的映射来思考。这一选择，连同我们表述中的种种技术细节，都是
由预期的几何应用（见 §17）所决定的。

\begin{lemma}\label{thm:quasi-representative-criterion}
(i) 取一对 $(c,k)$：
\begin{equation*}
\begin{aligned}
&c\in\hom_{\scr{A}}(Y_1,Y_2),\\
&k\colon\hom_{\scr{A}}(Y_0,Y_1)\longrightarrow\hom_{\scr{A}}(Y_0,Y_2)[-1],
\end{aligned}
\end{equation*}
它们满足
\begin{equation}\label{eq:pair-cocycle-equations}
\begin{aligned}
&\mu^1_{\scr{A}}(c)=0,\\
&\mu^1_{\scr{A}}\bigl(k(c_0)\bigr)-k\bigl(\mu^1_{\scr{A}}(c_0)\bigr)
+\mu^2_{\scr{A}}(c,c_0)=0
\quad\text{对一切 }c_0\in\hom_{\scr{A}}(Y_0,Y_1).
\end{aligned}
\end{equation}
这就确定了一个\trModule{}\trHomomorphism{}%
$t\colon\scr{T}_{Y_0}(\scr{Y}_1)\to\scr{Y}_2$，它具有性质
$[\mu^2_{\scr{Q}}(t,\iota)]=[\scr{J}^1(c)]\in\Hom_{H(\scr{Q})}(\scr{Y}_1,\scr{Y}_2)$。

(ii) 给定如 (i) 中的一对 $(c,k)$，以及任意的
\begin{equation*}
\begin{aligned}
&\gamma\in\hom_{\scr{A}}(Y_1,Y_2)[-1],\\
&\kappa\colon\hom_{\scr{A}}(Y_0,Y_1)\longrightarrow\hom_{\scr{A}}(Y_0,Y_2)[-2],
\end{aligned}
\end{equation*}
令
\begin{equation*}
\begin{aligned}
&\tilde{c}=c+\mu^1_{\scr{A}}(\gamma),\\
&\tilde{k}(c_0)=k(c_0)+(-1)^{|c_0|-1}\bigl(\mu^1_{\scr{A}}\bigl(\kappa(c_0)\bigr)
-\kappa\bigl(\mu^1_{\scr{A}}(c_0)\bigr)+\mu^2_{\scr{A}}(\gamma,c_0)\bigr),
\end{aligned}
\end{equation*}
便得到\eqref{eq:pair-cocycle-equations} 的一个新解，而与之相伴的\trHomomorphism{}%
$\tilde{t}$ 与原来的 $t$ \trCohomology。

(iii) 设 $(c,k)$ 如 (i) 中所述。此外，再假设对每个 $X\in\Ob\scr{A}$ 都给定了
映射
\begin{equation*}
\begin{aligned}
&d\colon\hom_{\scr{A}}(X,Y_1)\longrightarrow\hom_{\scr{A}}(X,Y_2),\\
&l\colon\hom_{\scr{A}}(Y_0,Y_1)\otimes\hom_{\scr{A}}(X,Y_0)
\longrightarrow\hom_{\scr{A}}(X,Y_2)[-1],\\
&\xi\colon\hom_{\scr{A}}(X,Y_1)\longrightarrow\hom_{\scr{A}}(X,Y_2)[-1],\\
&\psi\colon\hom_{\scr{A}}(Y_0,Y_1)\otimes\hom_{\scr{A}}(X,Y_0)
\longrightarrow\hom_{\scr{A}}(X,Y_2)[-2]
\end{aligned}
\end{equation*}
满足
\begin{equation*}
\medmuskip=1.5mu\relax\thickmuskip=3mu\relax
\begin{aligned}
&\mu^1_{\scr{A}}\bigl(d(c_1)\bigr)+d\bigl(\mu^1_{\scr{A}}(c_1)\bigr)=0,\\
&d\bigl(\mu^2_{\scr{A}}(c_0,b)\bigr)+\mu^1_{\scr{A}}\bigl(l(c_0,b)\bigr)
+(-1)^{|b|-1}l\bigl(\mu^1_{\scr{A}}(c_0),b\bigr)
+l\bigl(c_0,\mu^1_{\scr{A}}(b)\bigr)=0,\\
&\xi\bigl(\mu^1_{\scr{A}}(c_1)\bigr)-\mu^1_{\scr{A}}\bigl(\xi(c_1)\bigr)+d(c_1)
-\mu^2_{\scr{A}}(c,c_1)=0,\\
&(-1)^{|b|-1}\psi\bigl(\mu^1_{\scr{A}}(c_0),b\bigr)
+\psi\bigl(c_0,\mu^1_{\scr{A}}(b)\bigr)
-\mu^1_{\scr{A}}\bigl(\psi(c_0,b)\bigr)\\
&\quad+\xi\bigl(\mu^2_{\scr{A}}(c_0,b)\bigr)+l(c_0,b)
-\mu^2_{\scr{A}}\bigl(k(c_0),b\bigr)-\mu^3_{\scr{A}}(c,c_0,b)=0
\end{aligned}
\end{equation*}
（这里 $c_0$、$b$、$c_1$ 任取），并且使得下述\trChainComplex{}是\trAcyclic{}的：
\begin{equation}\label{eq:acyclic-cone-complex}
\begin{aligned}
D&=\hom_{\scr{A}}(Y_0,Y_1)\otimes\hom_{\scr{A}}(X,Y_0)[2]
\oplus\hom_{\scr{A}}(X,Y_1)[1]\oplus\hom_{\scr{A}}(X,Y_2),\\
\partial(c_0\otimes b,c_1,c_2)
&=\bigl((-1)^{|b|-1}\mu^1_{\scr{A}}(c_0)\otimes b
+c_0\otimes\mu^1_{\scr{A}}(b),\\
&\quad\mu^1_{\scr{A}}(c_1)+\mu^2_{\scr{A}}(c_0,b),\
\mu^1_{\scr{A}}(c_2)+d(c_1)+l(c_0,b)\bigr).
\end{aligned}
\end{equation}
那么，由 $(c,k)$ 定义的\trHomomorphism{} $t$ 诱导 $H^0(\scr{Q})$ 中的一个%
\trIsomorphism。事实上，只需对那样一组\trObject{}验证这一条件：它们构成%
\trFull{}子\trCategory{} $\tilde{\scr{A}}\subset\scr{A}$，而后者与 $\scr{A}$ %
\trQuasiequivalence；随后它对其余一切\trObject{}便会自动成立。
\end{lemma}

\begin{proof}
(i) 直接计算表明，如果我们如下定义 $t$，它便具有全部所需的性质：
\begin{equation*}
\begin{aligned}
t^d\bigl((c_0\otimes b,c_1),&a_{d-1},\ldots,a_1\bigr)\\
&=\mu^{d+1}_{\scr{A}}(c,c_1,a_{d-1},\ldots,a_1)
+\mu^{d+1}_{\scr{A}}\bigl(k(c_0),b,a_{d-1},\ldots,a_1\bigr)\\
&\quad+\mu^{d+2}_{\scr{A}}(c,c_0,b,a_{d-1},\ldots,a_1).
\end{aligned}
\end{equation*}

(ii) 完全照 (i) 那样做，只是把 $(c,k)$ 换成 $(\gamma,\kappa)$，便可以定义一个
$u\in\hom^{-1}_{\scr{Q}}(\scr{T}_{Y_0}\scr{Y}_1,\scr{Y}_2)$，它满足
$\mu^1_{\scr{Q}}(u)=\tilde{t}-t$。

(iii) 要使 $[t]$ 成为一个\trIsomorphism，我们需要 $t^1$ 的各个\trMappingCone{}%
都是\trAcyclic{}的。具体说来，这些\trCone{}就是如下的\trChainComplex：
\begin{equation*}
\medmuskip=1.5mu\relax\thickmuskip=3mu\relax
\begin{aligned}
\widetilde{D}&=\hom_{\scr{A}}(Y_0,Y_1)\otimes\hom_{\scr{A}}(X,Y_0)[2]
\oplus\hom_{\scr{A}}(X,Y_1)[1]\oplus\hom_{\scr{A}}(X,Y_2),\\
\widetilde{\partial}(c_0\otimes b,c_1,c_2)
&=\bigl((-1)^{|b|-1}\mu^1_{\scr{A}}(c_0)\otimes b
+c_0\otimes\mu^1_{\scr{A}}(b),\
\mu^1_{\scr{A}}(c_1)+\mu^2_{\scr{A}}(c_0,b),\\
&\quad\mu^1_{\scr{A}}(c_2)+\mu^2_{\scr{A}}(c,c_1)
+\mu^2_{\scr{A}}\bigl(k(c_0),b\bigr)+\mu^3_{\scr{A}}(c,c_0,b)\bigr).
\end{aligned}
\end{equation*}
利用 $(\xi,\psi)$ 可以构造一个\trIsomorphism{} $D\to\widetilde{D}$，即
\begin{equation*}
(c_0\otimes b,c_1,c_2)\longmapsto
\bigl(c_0\otimes b,c_1,c_2+\xi(c_1)+\psi(c_0\otimes b)\bigr).
\end{equation*}
于是\eqref{eq:acyclic-cone-complex} 的\trAcyclic{}性便蕴含了所要的命题。最后那句话是
显然的，因为限制\trFunctor{} $\scr{Q}=\mathit{mod}(\scr{A})\to
\mathit{mod}(\tilde{\scr{A}})$ 是一个\trQuasiequivalence
（\cref{thm:restriction-quasi-equivalence}）。
\end{proof}

\begin{lemma}\label{thm:third-morphism-realization}
取如\cref{thm:quasi-representative-criterion}(i) 中的一对 $(c,k)$。再假设此外
还有
\begin{equation*}
\medmuskip=1.5mu\relax\thickmuskip=3mu\relax
\begin{aligned}
g&=\sum_j g_2^j\otimes g_1^j\in
\bigl(\hom_{\scr{A}}(Y_0,Y_1)\otimes\hom_{\scr{A}}(Y_2,Y_0)\bigr)^{1},\\
p&=\sum_j p_2^j\otimes p_1^j\in
\bigl(\hom_{\scr{A}}(Y_0,Y_1)\otimes\hom_{\scr{A}}(Y_1,Y_0)\bigr)^{0},\\
q&=\sum_j q_2^j\otimes q_1^j\in\bigl(\hom_{\scr{A}}(Y_0,Y_1)\otimes
\Hom_{\KK}\bigl(\hom_{\scr{A}}(Y_0,Y_1),\hom_{\scr{A}}(Y_0,Y_0)\bigr)\bigr)^{-1},
\end{aligned}
\end{equation*}
满足
\begin{equation*}
\begin{aligned}
&\sum_j(-1)^{|g_1^j|-1}\mu^1_{\scr{A}}(g_2^j)\otimes g_1^j
+g_2^j\otimes\mu^1_{\scr{A}}(g_1^j)=0,\\
&\sum_j g_2^j\otimes\mu^2_{\scr{A}}(g_1^j,c)
=\sum_j(-1)^{|p_1^j|-1}\mu^1_{\scr{A}}(p_2^j)\otimes p_1^j
+\sum_j p_2^j\otimes\mu^1_{\scr{A}}(p_1^j),\\
&\sum_j g_2^j\otimes\bigl(\mu^2_{\scr{A}}\bigl(g_1^j,k(\cdot)\bigr)
+\mu^3_{\scr{A}}(g_1^j,c,\cdot)\bigr)\\
&\quad=\sum_j q_2^j\otimes\mu^1_{\scr{A}}\bigl(q_1^j(\cdot)\bigr)
-\sum_j q_2^j\otimes q_1^j\bigl(\mu^1_{\scr{A}}(\cdot)\bigr)\\
&\qquad+\sum_j(-1)^{|q_1^j|+|\cdot|-1}\mu^1_{\scr{A}}(q_2^j)\otimes q_1^j(\cdot)
+\sum_j p_2^j\otimes\mu^2_{\scr{A}}(p_1^j,\cdot)\\
&\qquad+\sum_j b^j\otimes\bigl(e_{Y_0}\otimes b^{j,\vee}\bigr),
\end{aligned}
\end{equation*}
其中：$e_{Y_0}$ 是 $Y_0$ 的\trIdentityMorphism{}的一个代表元；$b^j$ 是
$\hom_{\scr{A}}(Y_0,Y_1)$ 的一组\trBasis；$b^{j,\vee}$ 是相应的对偶\trBasis；
而 $(e_{Y_0}\otimes b^{j,\vee})\in\Hom_{\KK}\bigl(\hom_{\scr{A}}(Y_0,Y_1),
\hom_{\scr{A}}(Y_0,Y_0)\bigr)$ 是把 $b^j$ 送到 $e_{Y_0}$、而把其余\trBasis{}元素
全部变成零的那个\trLinear{}映射。那么，$g$ 给出一个\trDegree{}为一的%
\trModule{}\trHomomorphism{} $u\colon\scr{Y}_2\to
\hom_{\scr{A}}(Y_0,Y_1)\otimes\scr{Y}_0$，使得 $[\mu^2_{\scr{Q}}(u,t)]=[\pi]$
是来自\eqref{eq:twist-exact-triangle} 的那个映射。
\end{lemma}

\begin{proof}
定义预\trModule\trHomomorphism{}
\begin{equation*}
\begin{aligned}
&u\colon\scr{Y}_2\longrightarrow\hom_{\scr{A}}(Y_0,Y_1)\otimes\scr{Y}_0,\\
&u^d(b,a_{d-1},\ldots,a_1)
=\sum_j g_2^j\otimes\mu^{d+1}_{\scr{A}}(g_1^j,b,a_{d-1},\ldots,a_1),
\end{aligned}
\end{equation*}
以及
\begin{equation*}
\begin{aligned}
&v\colon\scr{T}_{Y_0}(\scr{Y}_1)\longrightarrow
\hom_{\scr{A}}(Y_0,Y_1)\otimes\scr{Y}_0,\\
&v^d\bigl((c_0\otimes b,c_1),a_{d-1},\ldots,a_1\bigr)
=\sum_j p_2^j\otimes\mu^{d+1}_{\scr{A}}(p_1^j,c_1,a_{d-1},\ldots,a_1)\\
&\quad+\sum_j q_2^j\otimes\mu^{d+1}_{\scr{A}}\bigl(q_1^j(c_0),b,a_{d-1},
\ldots,a_1\bigr)\\
&\quad+\sum_j p_2^j\otimes\mu^{d+2}_{\scr{A}}(p_1^j,c_0,b,a_{d-1},\ldots,a_1)\\
&\quad-\sum_j g_2^j\otimes\mu^{d+2}_{\scr{A}}(g_1^j,c,c_1,a_{d-1},\ldots,a_1)\\
&\quad-\sum_j g_2^j\otimes\mu^{d+2}_{\scr{A}}\bigl(g_1^j,k(c_0),b,a_{d-1},
\ldots,a_1\bigr)\\
&\quad-\sum_j g_2^j\otimes\mu^{d+3}_{\scr{A}}(g_1^j,c,c_0,b,a_{d-1},\ldots,a_1)
\end{aligned}
\end{equation*}
二者\trDegree{}分别为 $0$ 与 $-1$\footnote{译注：此处原书印作 “of degree 0 and
$-1$, respectively”；疑应为 “degree 1 and 0”（引理陈述中说 $u$ 的\trDegree{}为
一，且仅当 $\mu^2_{\scr{Q}}(u,t)$ 与 $\mu^1_{\scr{Q}}(v)$ 同为一次时才谈得上与
一次的 $\pi$ \trCohomology）；两处按数学内容当以此处为误，译文照原书排印，存疑备案。
——译者注}。利用这些假设可以算出 $\mu^1_{\scr{Q}}(u)=0$，而
$w=\mu^2_{\scr{Q}}(u,t)-\mu^1_{\scr{Q}}(v)$ 由
\begin{equation*}
w^d\bigl((c_0\otimes b,c_1),a_{d-1},\ldots,a_1\bigr)
=(-1)^{|c_0|}c_0\otimes\mu^{d+1}_{\scr{A}}(e_{Y_0},b,a_{d-1},\ldots,a_1).
\end{equation*}
给出。用与\cref{thm:cone-alternative-representatives}相同的论证可知，它与
$\pi$ \trCohomology。
\end{proof}

\begin{remark}\label{thm:exactness-from-conditions}
设我们处于\cref{thm:quasi-representative-criterion}的 (i) 与 (iii) 所描述的
情形；再假设 $\scr{A}$ 对有限\trDirectSum{}与\trSuspension{}封闭，并且
$\Hom_{H(\scr{A})}(Y_0,Y_1)$ 是有限维的，从而 $\Hom_{H(\scr{A})}(Y_0,Y_1)\otimes
Y_0$ 有意义。于是按构造，我们在 $H(\scr{A})$ 中有如下的\trExactTriangle
\begin{equation}\label{eq:homology-evaluation-triangle}
\begin{tikzcd}
\Hom_{H(\scr{A})}(Y_0,Y_1)\otimes Y_0 \ar[rr,"{[f]}"] && Y_1
\ar[dl,"{[c]}"] \\
& Y_2. \ar[ul,"{[1]}"] &
\end{tikzcd}
\end{equation}
从这一视角看，扭转过程退居幕后：我们手里只是一份条件清单（由%
\cref{thm:quasi-representative-criterion}给出），凭借它便可断定形如
\eqref{eq:homology-evaluation-triangle} 的三角形是\trExact{}的。不过要注意，
三角形中的第三个\trMorphism{}只有间接的刻画：它在\trYonedaEmbedding{}下的像是
$[\pi]\circ[t]^{-1}$。\cref{thm:third-morphism-realization}正是在这里变得有用。
设我们拥有该处所述的附加数据，并考虑 $H(\scr{A})$ 中的\trMorphism
\begin{equation*}
[g]=\sum_j[g_2^j]\otimes[g_1^j]\colon Y_2\longrightarrow
\Hom_{H(\scr{A})}(Y_0,Y_1)\otimes Y_0
\end{equation*}
在\trYonedaEmbedding{}下，它被映到\trModule\trHomomorphism{} $[u]$，而按构造有
$[u]\circ[t]=[\pi]$。因此，\eqref{eq:homology-evaluation-triangle} 中的第三个
映射恰好就是 $[g]$。
\end{remark}
